\chapter{Att kontrollera ett påstående}
\label{app:verifiera-mangder}

Det här materialet gör, som allt undervisningsmaterial, en rad
påståenden.  De flesta kan du pröva själv genom att köra programmen ---
det är hela poängen med att koden i materialet är samma kod som körs.  Men
några påståenden går inte att pröva på det sättet:
\begin{enumerate}
  \item att medlemskapstestet i en mängd i genomsnitt kostar lika mycket
    oavsett hur många element det finns, medan samma test i en lista
    kostar mer ju längre listan är (\cref{sec:mangdval});
  \item att en \mintinline{python}{deque} är ett bättre val än en lista
    för en kö, därför att listan måste flytta alla element när det
    första tas bort (\cref{sec:stackar});
  \item att en mängd varken har dubbletter eller ordning
    (\cref{sec:mangder})%
  % Claims made only in the instructor notes print only in the teacher
  % edition.
  \ifbool{ltnote}{;
  \item att lösningen ska hållas enkel --- principen
    \foreignlanguage{english}{KISS} (\cref{sec:valja-mangder});
  \item de didaktiska påståenden som anteckningarna i marginalen gör om
    hur människor lär sig.}{.}
\end{enumerate}
Hur vet man om sådant stämmer?  Samma fråga möter dig varje gång du läser
en lärobok, en webbsida eller ett svar från en språkmodell --- och varje
gång du själv ska skriva en rapport.
\Cref{tab:mangder-pastaenden} visar var vart och ett av materialets
påståenden är belagt.  Lägg särskilt märke till de två första: en siffra
om hur snabb en operation är ser exakt ut som ett faktum, men gäller bara
under förutsättningar som måste anges.

\begin{table}[!htb]
  \centering
  \caption{Modulens påståenden, som frågor, och var svaret finns.}
  \label{tab:mangder-pastaenden}
  \begin{tabular}{@{}p{0.46\linewidth}p{0.44\linewidth}@{}}
    \toprule
    Fråga & Var det är belagt \\
    \midrule
    Kostar ett medlemskapstest i en mängd lika mycket oavsett antalet
    element, till skillnad från samma test i en lista?
      & Primärkälla: Pythons egen sammanställning av operationernas
        kostnad, läst i fulltext 2026-09-03.  Ja --- i genomsnitt, i
        CPython; i värsta fall nej. \\
    Är en \mintinline{python}{deque} ett bättre val än en lista för en
    kö?
      & Primärkälla: Pythons dokumentation för modulen
        \mintinline{python}{collections}, läst i fulltext 2026-09-03. \\
    Saknar en mängd dubbletter och ordning?
      & Primärkälla: Pythons egen dokumentation, läst i fulltext
        2026-09-03; dessutom visar programmet i \cref{sec:mangder} det. \\
    % Rows for claims made only in the instructor notes print only in
    % the teacher edition.
    \ifbool{ltnote}{%
    Ska lösningen hållas enkel (\foreignlanguage{english}{KISS})?
      & Redan belagt i föreläsningen \emph{Funktioner}, bilaga~C: att
        enkelhet är ett etablerat designvärde, medan den vanliga
        upphovsuppgiften inte går att bekräfta. \\
    Måste kontrast upplevas samtidigt för att en aspekt ska kunna
    urskiljas?
      & Primärkälla: Marton, \emph{Necessary Conditions of Learning},
        s.~45--47. \\
    }{}%
    \bottomrule
  \end{tabular}
\end{table}

Arbetsgången är densamma i alla kursens föreläsningar och beskrivs i
föreläsningen \emph{Algoritmiskt tänkande}, bilaga~A: hur ett påstående
görs till en fråga som kan besvaras, hur källan väljs efter frågan och
söks fram ämnesdrivet och tvåsidigt, med en motsökning efter
invändningar, hur fulltexten läses och spåret skrivs ner, och hur
slutsatsen dras med sina förbehåll.  \Cref{sec:mangder-pastaenden}
säger var den här föreläsningens påståenden är belagda.

\section{Materialets påståenden och var de är belagda}
\label{sec:mangder-pastaenden}

Inget av den här modulens påståenden krävde en egen litteratursökning.
Frågor om vad Python gör besvaras av Pythons egen dokumentation, som är
primärkällan: vad en mängd är och att den saknar
dubbletter\autocite{PythonDataStructures}, hur mycket varje operation
kostar\autocite{PythonTimeComplexity}, och varför
\mintinline{python}{collections.deque} finns\autocite{PythonCollections}.
Dokumentationen för \mintinline{python}{collections} säger det rakt
ut: listor \enquote{are optimized for fast
fixed-length operations and incur \(O(n)\) memory movement costs for
\mintinline{python}{pop(0)}}.  Alternativet --- att ta tid på
operationerna --- hade mätt den här datorn och inte datastrukturen.
Komplexitetssidan säger själv att den beskriver \enquote{current
CPython} och att siffrorna är genomsnitt, och det förbehållet står
därför också i brödtexten.

\ifbool{ltnote}{%
Principen att välja den enklaste behållare som klarar de operationer man
behöver (\foreignlanguage{english}{KISS}) namnges i lärarnoterna; det
vetenskapliga
underlaget för \foreignlanguage{english}{KISS} --- att enkelhet är ett
etablerat designvärde, medan den vanliga upphovsuppgiften inte går att
bekräfta --- finns i föreläsningen \emph{Funktioner}, bilaga~C.
Att kontrast måste upplevas samtidigt för
att en aspekt ska kunna urskiljas är hämtat från Martons
\emph{Necessary Conditions of Learning}\autocite[45--47]{Marton2015}.
}{}
